% Part: first-order-logic
% Chapter: introduction
% Section: formulas

\documentclass[../../../include/open-logic-section]{subfiles}

\begin{document}

\olfileid{fol}{int}{fml}

\section{\usetoken{P}{formula}}

Iki cara kang bakal kita gunakake kanggo nemtokake kanthi trep
!!{sentence}s ing logika sepisanan lan ngrampungi prekara kang tuwuh
saka panganggone !!{variable}s. Sepisan kita netepake sawijining
himpunan ekspresi kang \emph{beda}: !!{formula}s. Sawise iku rampung,
kita bisa nimbang kalungguhane !!{variable}s ing formula-formula
kasebut---lan adhedhasar sawetara gagasan liyane, yaiku !!{variable}s
kang ``bebas'' lan kang ``kaiket,'' kita bisa nemtokake apa iku
!!a{sentence} (yaiku !!a{formula} tanpa !!{variable}s bebas). Iki
ditindakake ora mung amarga ndadekake definisi ``!!{sentence}'' luwih
gampang ditata, nanging uga amarga bakal wigati banget tumrap cara kita
nemtokake gagasan semantis ngenani relasi nyukupi.

Ayo kita netepake ``!!{formula}'' kanggo basa logika sepisanan kang
prasaja, kang mung ngemot siji !!{predicate}~$\Obj P$ lan siji
!!{constant}~$\Obj a$, sarta mung simbol logis $\lnot$, $\land$,
lan~$\lexists$. Definisi kita kang jangkep bakal luwih umum:
kita bakal marengake !!{predicate}s lan !!{constant}s kang cacahe tanpa
wates. Malah, kita uga bakal nimbang !!{function}s kang bisa digabungake
karo !!{constant}s lan !!{variable}s kanggo mbentuk ``term.'' Sauntara
iki, $\Obj a$ lan variabel-variabel mau bakal dadi siji-sijine term
kita. Nanging, kita pancen mbutuhake !!{variable}s kang cacahe tanpa
wates. Kanthi resmi kita bakal nggunakake simbol $\Obj v_0$, $\Obj
v_1$, \dots, minangka variabel.

\begin{defn}
Himpunan \emph{!!{formula}s}~$\Frm$ ditegesi kaya mangkene:
\begin{enumerate}
\item\ollabel{fmls-atom} $\Atom{\Obj P}{\Obj a}$ lan $\Atom{\Obj
  P}{\Obj v_i}$ yaiku !!{formula}s ($i \in \Nat$).

\tagitem{prvNot}{\ollabel{fmls-not}Yen $!A$ iku !!a{formula}, mula $\lnot !A$
  yaiku !!{formula}.}{}

\tagitem{prvAnd}{Yen $!A$ lan $!B$ iku !!{formula}s, mula $(!A \land
  !B)$ iku !!a{formula}.}{}

\tagitem{prvEx}{\ollabel{fmls-ex}Yen $!A$ iku !!a{formula} lan $x$ iku !!a{variable},
  mula $\lexists[x][!A]$ iku !!a{formula}.}{}

\tagitem{limitClause}{\ollabel{fmls-limit}Ora ana liyane kang dadi !!a{formula}.}{}
\end{enumerate}
\end{defn}

\olref{fmls-atom} ngandhani kita yen $\Atom{\Obj P}{\Obj a}$ lan
$\Atom{\Obj P}{\Obj v_i}$ iku !!{formula}s, kanggo saben $i \in
\Nat$. Iki diarani !!{formula}s \emph{atom}. Formula-formula iki
menehi titik wiwitan. Klausa-klausa liyane menehi cara kanggo mbentuk
!!{formula}s anyar saka formula kang wis kita bentuk. Dadi, tuladhane,
miturut \olref{fmls-not}, kita oleh asil yen $\lnot \Atom{\Obj P}{\Obj
v_2}$ iku !!a{formula}, amarga $\Atom{\Obj P}{\Obj v_2}$ wis dadi
!!a{formula} miturut \olref{fmls-atom}. Banjur, miturut
\olref{fmls-ex}, kita oleh asil yen $\lexists[\Obj v_2][\lnot
\Atom{\Obj P}{\Obj v_2}]$ iku !!{formula} liyane, lan sabanjure.
\olref{fmls-limit} ngandhani kita yen \emph{mung} rerangken kang bisa
kita bentuk nganggo cara iki kang dianggep !!{formula}s. Mligine,
$\lexists[\Obj v_0][\Atom{\Obj P}{\Obj a}]$ lan $\lexists[\Obj
  v_0][\lexists[\Obj v_0][\Atom{\Obj P}{\Obj a}]]$ \emph{pancen}
dianggep !!{formula}s, dene $(\lnot \Atom{\Obj P}{\Obj a})$ ora,
amarga kurung njaba kasebut keluwih.

Cara nemtokake !!{formula}s kaya iki diarani \emph{definisi
induktif}, lan ndadekake kita bisa mbuktekake prekara ngenani
!!{formula}s nganggo salah sawijining ragam bukti kanthi induksi kang
diarani \emph{induksi struktural}. Bab iki dirembug kanthi umum ing
\olref[mth][ind][idf]{sec} lan \olref[mth][ind][sti]{sec}, kang kudu
sampeyan sinau maneh sadurunge nyinaoni bukti-bukti sabanjure. Gagasan
dhasare yaiku yen sampeyan arep menehi bukti yen sawijining prekara
bener tumrap kabeh !!{formula}s, sepisan sampeyan nuduhake yen prekara
iku bener tumrap !!{formula}s atom, banjur nuduhake yen \emph{yen}
prekara iku bener tumrap sembarang !!{formula}~$!A$ (lan~$!B$), prekara
iku \emph{uga} bener tumrap $\lnot !A$, $(!A \land !B)$, lan
$\lexists[x][!A]$. Tuladhane, iki mbuktekake yen prekara kasebut bener
tumrap $\lexists[\Obj v_2][\lnot \Atom{\Obj P}{\Obj v_2}]$: saka
perangan kapisan sampeyan ngerti yen prekara iku bener tumrap
!!{formula} atom~$\Atom{\Obj P}{\Obj v_2}$. Banjur, saka perangan
kapindho sampeyan oleh asil yen prekara iku bener tumrap $\lnot
\Atom{\Obj P}{\Obj v_2}$, lan banjur oleh maneh yen prekara iku bener
tumrap $\lexists[\Obj v_2][\lnot \Atom{\Obj P}{\Obj v_2}]$ dhewe.
Amarga kabeh !!{formula}s diasilake kanthi induktif saka !!{formula}s
atom, cara iki lumaku tumrap formula endi wae.

\end{document}
